% Generated draft: requires mathematical and source review.
\documentclass[11pt]{article}
\usepackage[margin=25mm]{geometry}
\usepackage{fontspec}
\setmainfont{Latin Modern Roman}
\usepackage{amsmath,amssymb,mathtools,unicode-math}
\setmathfont{Latin Modern Math}
\usepackage{enumitem,needspace,xurl,hyperref}
\setlength{\parindent}{0pt}
\setlength{\parskip}{5pt}
\setlength{\emergencystretch}{3em}
\newcommand{\N}{\mathbb{N}}
\newcommand{\Z}{\mathbb{Z}}
\newcommand{\Q}{\mathbb{Q}}
\newcommand{\R}{\mathbb{R}}
\newcommand{\C}{\mathbb{C}}
\newcommand{\F}{\mathbb{F}}
\newcommand{\A}{\mathbb{A}}
\newcommand{\PP}{\mathbb{P}}
\newcommand{\E}{\mathbb{E}}
\newcommand{\eps}{\varepsilon}
\newcommand{\dd}{\,\mathrm{d}}
\begin{document}
\section*{UnsolvedMath ID: 30001136}
\textbf{Source problem code:} catalogue entry
\subsection{Definitions and mathematical statement}
Let \(K_4\) be the complete graph on vertices \(1,2,3,4\). A planar realization is a tuple \(P=(p_1,p_2,p_3,p_4)\) of distinct points of \(\R^2\), with an edge represented by the segment joining its endpoints. The supplied record requires general position but does not define it. The usual condition that no three points are collinear is a possible interpretation, requiring confirmation. It is also unspecified whether crossings are allowed or whether ``embedding'' requires a crossing-free drawing.

A self-stress of \(P\) is a collection \(\omega=(\omega_{ij})_{1\leq i<j\leq4}\), extended by \(\omega_{ji}=\omega_{ij}\), satisfying
\[
\sum_{j\ne i}\omega_{ij}(p_j-p_i)=0
\qquad (i=1,2,3,4).
\]
For an affinely spanning realization of \(K_4\), the self-stress space is one-dimensional. Consequently, specifying an associated nonzero self-stress requires a choice of sign; a strict positivity condition cannot be interpreted independently of that choice.

Write \(\mathcal G\) for the intended class of general-position realizations, \(\omega(P)\) for the intended selected self-stress, and \(\mathcal S\) for the intended domain of segments. These are placeholders for source conventions that the supplied record does not fully specify. If \(f:\mathcal S\to\R\) is a single function evaluated on the edges of every \(P\in\mathcal G\), put
\[
f_{ij}(P)=f(\overline{p_ip_j}),\qquad
L_P(f)=\sum_{1\leq i<j\leq4}\omega_{ij}(P)f_{ij}(P).
\]
\textbf{Requested characterization, conditional on the missing conventions.} Characterize all admissible functions \(f\) satisfying
\[
\forall P\in\mathcal G,\qquad L_P(f)>0.
\]
The record does not specify the admissible function class, any regularity assumptions, or whether this global segment-function interpretation is intended instead of a configuration-dependent assignment.

For precision about the proposed equivalence, one conventional rigidity map is
\[
R_P:(\R^2)^4\longrightarrow\R^6,\qquad
(R_Pu)_{ij}=(p_i-p_j)\mathbin{\cdot}(u_i-u_j).
\]
Its range consists of edge assignments arising from vertex velocities, with this normalization of the map. The source does not specify how additions from these configuration-dependent ranges are to define a single admissible function on segments.

Under the conventional interpretation of brackets as signed doubled triangle area, define
\[
[xyz]=\det(y-x,z-x),\qquad
F_{a,b}(\overline{xy})=[axy][byx],
\]
where \(a,b\in\R^2\). The product is unchanged when \(x\) and \(y\) are interchanged. The record asks whether every function with the required positivity is obtained from this construction by positive scaling and addition of a function in the rigidity-map range. At an individual realization, the displayed candidate form would read
\[
\bigl(f_{ij}(P)\bigr)_{i<j}
=c\bigl([ap_ip_j][bp_jp_i]\bigr)_{i<j}+R_Pu,
\qquad c>0,\quad u\in(\R^2)^4.
\]
This expression records the proposed form, not a settled universal formulation: the record supplies neither restrictions on \(a,b\) nor the quantifiers determining whether \(a,b,c,u\) may depend on \(P\). The bracket convention also requires confirmation. Without these conventions and a stress sign selection, the requested classification is not fully specified.
\subsection{Short English statement}
For permanent UnsolvedMath ID \(30001136\), the question asks which assignments to planar segments have a strictly positive pairing with the selected self-stress of every admissible realization of \(K_4\), and whether the displayed area-product construction generates all such assignments under the stated operations. The stress sign, realization class, function domain, auxiliary-point restrictions, and equivalence quantifiers remain unresolved. This is a statement-only draft requiring human mathematical review.
\subsection{Sources}
\begin{itemize}
\item Local registry: supplied complete source record for permanent UnsolvedMath ID \(30001136\), titled ``Positive Self-Stress Functionals on Planar \(K_4\) Embeddings,'' in the Oberwolfach Reports Open Problems collection. Exact source problem code: \texttt{catalogue entry}. The record gives proposed year \(2009\) and names no proposer.
\item Supplied source citation: \emph{Discrete Differential Geometry} (2009), Oberwolfach Reports, pp.~75--144, DOI \texttt{10.4171/owr/2009/02}. \url{https://doi.org/10.4171/owr/2009/02}. The background describes these reports as documenting mathematical workshops and their open-problem discussions.
\item Supplied literature reference: Rote--Santos--Streinu, ``expansive motions.'' \url{https://arxiv.org/abs/math/0206027}.
\item Supplied literature reference: Rote, \emph{Open Problems in Discrete Differential Geometry}. \url{https://page.mi.fu-berlin.de/rote/Kram/OWR-DDG09-problems.pdf}.
\item Literature assessment dated exactly \texttt{2026-08-21}: the supplied curation labels the question open and states that the cited construction supplies examples but not an exhaustive classification. This is supplied dated curation derived from the OWR record, not an independently verified literature assessment; time-sensitive status and cited claims require independent verification.
\item The registry labels the statement \texttt{corrected\_verified} and supplies the verification text ``The positivity classification is complete; the displayed stress functional and candidate family are normalized.'' These are supplied metadata claims, not mathematical verification performed for this draft, and their meaning requires reconciliation with the open classification question and the missing conventions above.
\end{itemize}
\end{document}
